\begin{abstract}
  We introduce bindings that enable the convenient, efficient, and
  reliable use of software modules of \cgal{} (Computational Geometry
  Algorithm Library), which are written in C++, in code written in
  Python.
  %
  There are different ways to go about the implementation of such
  bindings. We present a short study that compares three main methods,
  which leads to the method of choice.
  % The implementation of practical algorithms and data structures in
  computational geometry presents tremendous difficulties, such as
  obtaining stable software despite the use of (inexact) floating
  point arithmetic, found in standard hardware, and obtaing complete
  software that handles all degenerate cases, which typically are in
  abundance. The code of \cgal{} extenssively uses function and class
  templates in order to handle these difficulties, which implies that
  the programmer has to make many choices that are resolved during
  compile time (of the C++ modules). While bindings take effect at run
  time (of the Python code), the type of the C++ objects that are
  bound, must be known when the bindings are generated. The types of
  the bound objects are instances (instantiated types) of C++ function
  and class templates. The number of object types that potentially can
  be bound enormous, and the generation of the bindings in advance
  practiaclly impossible. We present a system that dynamically
  generates bindings for desired object types according to user
  prescriptions, which enables the convenient use of any subset of
  bound object types concurrently.
\end{abstract}
